% !TEX root = ../../ctfp-print.tex

\lettrine[lhang=0.17]{自}{由构造是伴随函子的一个强大应用。}一个\newterm{自由函子（free functor）}被定义为某个\newterm{遗忘函子（forgetful functor）}的左伴随。遗忘函子通常是一个相对简单的函子，它会"忘记"某些结构。例如，许多有趣的范畴都是基于集合构建的。但范畴论中的对象虽然抽象了这些集合，却本身没有内部结构——它们没有元素。不过这些对象往往保留着集合的记忆，从某种意义上说存在一个映射（即函子），将给定范畴$\cat{C}$到$\Set$进行关联。与$\cat{C}$中某个对象对应的集合称为其\newterm{底层集合（underlying set）}。

幺半群就是具有底层集合的对象——即元素的集合。存在一个遗忘函子$U$，它将幺半群范畴$\cat{Mon}$映射到集合范畴，把幺半群映射为其底层集合，也将幺半群同态（即保持运算结构的映射）映射为集合间的函数。

我喜欢认为$\cat{Mon}$具有双重身份：一方面它是带有乘法和单位元的集合；另一方面它是一个由无特征对象构成的范畴，其唯一结构体现在对象间的态射中。每个保持乘法和单位元的集合函数都会产生$\cat{Mon}$中的一个态射。

需要注意：
\begin{itemize}
  \tightlist
  \item
        可能有多个不同的幺半群对应同一个底层集合；
  \item
        幺半群同态的数量少于（或至多等于）其底层集合间的函数数量。
\end{itemize}

作为遗忘函子$U$左伴随的函子$F$就是自由函子，它通过生成集构造自由幺半群。这种伴随关系源自我们之前讨论过的自由幺半群泛构造\footnote{参见第13章\hyperref[free-monoids]{自由幺半群}相关内容}。

\begin{figure}[H]
  \centering
  \includegraphics[width=0.6\textwidth]{images/forgetful.jpg}
  \caption{幺半群$m_1$和$m_2$具有相同底层集合。$m_2$与$m_3$的底层集合间函数数量多于它们之间的态射数量。}
\end{figure}

用态射集表示该伴随关系可写作：
$$\cat{Mon}(F x, m) \cong \Set(x, U m)$$
这个（关于$x$和$m$自然）的同构表明：

\begin{itemize}
  \tightlist
  \item
        对任意由$x$生成的自由幺半群$F x$到任意幺半群$m$的同态，都存在唯一的函数将生成集$x$嵌入$m$的底层集合中，即$\Set(x, U m)$中的函数。
  \item
        对任意将$x$嵌入某$m$底层集合的函数，都存在唯一的自由幺半群生成集到$m$的幺半群同态。（这就是我们在泛构造中称作$h$的那个态射）
\end{itemize}

\begin{figure}[H]
  \centering
  \includegraphics[width=0.8\textwidth]{images/freemonadjunction.jpg}
\end{figure}

直观上，$F x$是基于$x$能构造出的"最大"幺半群。如果我们能观察幺半群内部结构，会发现$\cat{Mon}(F x, m)$中的任何态射都通过识别某些元素的方式，将这个自由幺半群嵌入其他幺半群$m$中。特别地，它将$F x$的生成元（即$x$的元素）嵌入$m$。伴随性表明：右侧$\Set(x, U m)$给出的$x$嵌入方式唯一确定左侧的幺半群嵌入，反之亦然。

在Haskell中，列表数据结构就是一个自由幺半群（有一些例外情况：参见\urlref{http://comonad.com/reader/2015/free-monoids-in-haskell/}{Dan Doel的博客文章}）。类型\code{{[}a{]}}是一个自由幺半群，其中类型\code{a}代表生成集。例如，\code{{[}Char{]}}包含单位元空列表\code{{[}{]}}以及单元素列表如\code{{[}'a'{]}}, \code{{[}'b'{]}}（即自由幺半群的生成元）。其余元素通过"乘积"操作生成——这里两个列表的乘积就是简单拼接。拼接操作满足结合律和单位律（空列表为单位元）。由\code{Char}生成的自由幺半群本质上就是所有字符序列的集合，在Haskell中被称为\code{String}：

\src{snippet01}
（\code{type}定义了一个类型同义词——现有类型的另一个名称）

另一个有趣的例子是由单一生成元构造的自由幺半群，即单位元列表\code{{[}(){]}}。其元素包括\code{{[}{]}}、\code{{[}(){]}}、\code{{[}(), (){]}}等。每个这样的列表都可以用一个自然数（其长度）来描述。单位元列表不再携带其他信息。两个此类列表的拼接结果的新长度等于两者长度之和。容易看出\code{{[}(){]}}类型与自然数加法幺半群（含零）同构。以下是互逆的两个函数见证这种同构：

\src{snippet02}
为简化起见我使用了\code{Int}而非\code{Natural}，但核心思想一致。\code{replicate}函数创建一个长度为\code{n}的列表，所有元素填充为给定值（此处是单位元）。

\section{一些直观理解}

以下是一些非形式化的直观解释。这类论述虽不够严谨，但有助于建立直观认知。

理解自由/遗忘伴随函子时，记住函子和函数本质上都是信息丢失的过程会有帮助。函子可能会合并多个对象和态射，函数可能会将集合元素聚集在一起，且它们的像可能仅覆盖陪域的一部分。

在$\Set$范畴中，典型的态射集包含从最小信息丢失（如单射或可能的同构）到最大信息丢失（如将整个定义域坍缩为单个元素的常量函数）的各种函数。

我也倾向于认为任意范畴中的态射都具有信息丢失特性。这只是一个心智模型，但在思考伴随函子（尤其是涉及$\Set$的伴随）时非常有用。

形式上我们只能讨论可逆态射（同构）和不可逆态射。后者通常被认为是信息丢失的。还有单态射和满态射的概念推广了单射（不坍缩）和满射（覆盖整个陪域）函数的性质，但存在既是单态又是满态但仍然不可逆的态射。

在自由函子⊣遗忘函子伴随中，左侧是约束更强的范畴$\cat{C}$，右侧是约束更弱的范畴$\cat{D}$。$\cat{C}$中的态射"更少"因为它们必须保持额外结构。在$\cat{Mon}$情形，它们必须保持乘法和单位元。而$\cat{D}$中的态射不需要保持这么多结构，因此数量"更多"。

当我们应用遗忘函子$U$到$\cat{C}$中的对象$c$时，我们认为这是揭示了$c$的"内部结构"。事实上，当$\cat{D}$是$\Set$时，我们认为$U$定义了$c$的内部结构——即其底层集合。（在任意范畴中，我们不能讨论对象的内部细节，只能通过与其他对象的联系来认识，但这里只是非正式讨论）

当我们用$U$映射两个对象$c'$和$c$时，一般预期$\cat{C}(c', c)$的映射只会覆盖$\cat{D}(U c', U c)$的一个子集。这是因为$\cat{C}(c', c)$中的态射必须保持额外结构，而$\cat{D}$中的则不需要。

\begin{figure}[H]
  \centering
  \includegraphics[width=0.45\textwidth]{images/forgettingmorphisms.jpg}
\end{figure}

但由于伴随性被定义为特定态射集的同构，我们必须非常谨慎地选择$c'$。在伴随中，$c'$不是随意从$\cat{C}$中选取的，而是来自（可能更小的）自由函子$F$的像：
$$\cat{C}(F d, c) \cong \cat{D}(d, U c)$$
因此$F$的像必须由那些到任意$c$都有大量态射的对象组成。实际上，从$F d$到$c$的结构保持态射数量必须等于从$d$到$U c$的非结构保持态射数量。这意味着$F$的像必须由本质上无结构的对象（这样态射就不需要保持任何结构）组成。这种"无结构"对象被称为自由对象。

\begin{figure}[H]
  \centering
  \includegraphics[width=0.45\textwidth]{images/freeimage.jpg}
\end{figure}

在幺半群例子中，自由幺半群除了单位元和结合律生成的结构外没有其他结构。除此之外，所有乘法都会产生全新元素。

在自由幺半群中，$2 * 3$不是$6$，而是新元素${[}2, 3{]}$。由于${[}2, 3{]}$和$6$没有被识别为同一，从这个自由幺半群到任何其他幺半群$m$的态射可以分别映射它们。当然也可以将${[}2, 3{]}$和$6$映射到$m$的同一元素。或者在加法幺半群中识别${[}2, 3{]}$和$5$（它们的和）等等。不同的识别方式会产生不同的幺半群。

这引出了另一个有趣的直观：自由幺半群不是执行幺半群运算，而是积累传入的参数。它不像传统乘法那样直接计算$2$和$3$，而是将它们记录在列表中。这种方案的优势在于我们不必预先指定使用的幺半群运算。我们可以持续积累参数，直到最后才应用运算符。这时我们可以选择加法、乘法、模2加法等不同运算。自由幺半群将表达式的创建与其求值分离。当我们讨论代数时会再次看到这个思想。

这种直观可以推广到其他更复杂的自由构造。例如，我们可以先累积完整的表达式树再进行求值。这种方法的优势在于我们可以变换这些树以使求值更快或节省内存。例如在实现矩阵微积分时，这种延迟求值策略可以避免立即分配大量临时数组存储中间结果。

\section{挑战}

\begin{enumerate}
  \tightlist
  \item
        考虑由单元素集合作为生成元所构建的自由幺半群（free monoid）。
        证明：从该自由幺半群到任意幺半群 $ m $ 的态射（morphism），与从该单元素集合到 $ m $ 的底层集合（underlying set）的函数之间存在一一对应关系。
\end{enumerate}